\chapter{统一模仿学习策略与实验设计}\label{chap:policy-design}

\section{目标中心表示与动作块构造}

实验中所有学习策略均作为任务空间动作块规划器运行，而非端到端关节控制器。每个监督样本包含历史观测窗口、未来目标条件和专家动作块。历史长度记为$H_o$，主实验设$H_o=16$；动作块长度为$H_a=32$；未来条件长度为$H_c=50$。运行时每次规划32步动作块，但只执行前16步后重新规划，因此执行窗口为$H_{\mathrm{exec}}=16$。

目标中心表示把抓取问题转化为相对目标的接近与对齐——不需要在世界坐标中记录绝对末端轨迹。设立方体位置和姿态为$(p^c_t,q^c_t)$，末端执行器位置和姿态为$(p^e_t,q^e_t)$，$R(q)$为四元数对应旋转矩阵\citep{shoemake1985animating}，则相对末端位置和姿态为
\begin{align}
    \tilde{p}^e_t &= R(q^c_t)^\top(p^e_t-p^c_t),\\
    \tilde{q}^e_t &= (q^c_t)^{-1}\otimes q^e_t .
    \label{eq:target-centric-observation}
\end{align}
未来目标条件为7维序列，包含相对位置变化$\Delta p(3)$和相对四元数$\Delta q(4)$，给出未来50步的立方体运动预测。常速度外推模型根据当前观测生成该序列：位置以当前线速度线性外推，姿态以当前角速度逐步积分；外推结果再经目标中心坐标变换转为相对变化量。该构造不直接读取仿真中的未来真值——在真实部署中，只要能从感知系统估计出目标的当前速度，即可用相同的常速度模型在线生成未来条件。在随机姿态翻滚平移场景中，该外推只是近似（空间轴旋转时角速度方向本身在演化），但为策略给出了合理的短时运动参考。动作标签为8维序列，包含相对任务空间目标位置3维、相对目标四元数4维和夹爪值1维。主表示使用完整目标中心历史观测，其历史维度为14；紧凑目标中心表示、紧凑世界坐标表示和完整世界坐标表示供表示消融配置使用。

策略输出的相对动作块在运行时解码为世界坐标任务空间参考。若第$k$步预测动作为$(\hat{p}^{rel}_{t,k},\hat{q}^{rel}_{t,k},\hat{g}_{t,k})$，则对应世界坐标参考为
\begin{align}
    \hat{p}^{r}_{t+k} &= p^c_{t+k}+R(q^c_{t+k})\hat{p}^{rel}_{t,k},\\
    \hat{q}^{r}_{t+k} &= q^c_{t+k}\otimes \hat{q}^{rel}_{t,k}.
    \label{eq:relative-action-decode}
\end{align}
连续夹爪通道参与训练监督，但闭环执行中的夹爪命令由状态机决定，以避免连续回归在开闭突变附近产生不稳定拖尾。

\section{模仿学习策略}

本文比较四类共享输入输出接口的模仿学习策略：行为克隆（BC, Behavior Cloning）、循环行为克隆（BC-RNN, Recurrent Behavior Cloning）、动作分块变换器（ACT, Action Chunking Transformer）和扩散策略（DP, Diffusion Policy）。四类策略的输入均为历史观测窗口$O_t$和未来目标条件$C_t$，输出均为长度为$H_a=32$的任务空间动作块$\hat{A}_t$。因此，本节中模型的差异主要源于条件编码方式、动作块生成主干和训练目标，而非不同的观测、动作定义或跟踪执行器。

\subsection{统一策略接口与辅助预测头}

每个策略输出包含三部分：连续动作块、闭合预测头和阶段预测头。连续动作块为$32\times 8$序列，其中位置、姿态通道在运行时解码为任务空间末端参考；连续夹爪通道参与监督训练，但闭环执行中的最终夹爪命令由状态机下达。闭合预测头预测夹爪闭合事件，触发从接近或对齐到闭合的切换；阶段预测头预测五类专家阶段，约束状态机单调推进。这一设计让四类学习模型均担当高层任务空间动作块规划器，而非直接输出关节命令。

图~\ref{fig:policy-architectures}展示四类策略在统一输入输出接口下的结构差异；表~\ref{tab:model-architecture-summary}列出各主模型的架构摘要。DP其余容量变体（1.5M、6M、18M）见附录。

\begin{figure}[!htbp]
    \centering
    \includegraphics[width=0.95\textwidth]{fig_policy_architectures}
    \caption{统一动作块策略接口下的四类模仿学习模型。BC、BC-RNN、ACT和DP共享历史观测窗口、未来目标条件、32步任务空间动作块、闭合预测头和阶段预测头；差异主要源于条件编码方式、动作块生成主干和训练目标。}
    \label{fig:policy-architectures}
\end{figure}

\begin{table}[!htbp]
    \caption{模仿学习策略的主要架构设置。}
    \label{tab:model-architecture-summary}
    \centering
    \small
    \renewcommand{\arraystretch}{1.15}
    \begin{tabularx}{\textwidth}{>{\raggedright\arraybackslash}p{0.13\textwidth}>{\raggedright\arraybackslash}p{0.30\textwidth}Y}
        \hline
        策略 & 核心结构 & 主要设置 \\
        \hline
        BC & 历史因果1D-CNN、未来条件1D-CNN和MLP动作块回归头 & 历史通道$[64,128,192]$，未来通道$[64,128]$，条件维度192，MLP隐藏维度$[480,480,480]$ \\
        BC-RNN & GRU历史编码器、未来条件1D-CNN和GRU动作块解码器 & 条件维度192，2层GRU，隐藏维度218，约1.5M参数 \\
        ACT & 变换器编码器--解码器和CVAE后验编码器 & 变换器维度132，4头注意力，2层编码器、2层解码器，FFN维度544，隐变量维度32 \\
        DP 12M & 条件1D U-Net动作块去噪器 & 条件维度160，基础通道128，通道倍率$[1,2,4]$，100步训练扩散，10步DDIM推理 \\
        \hline
    \end{tabularx}
\end{table}

\subsection{确定性动作块策略：BC与BC-RNN}

BC是最简单直接的动作块回归基线。它以因果一维卷积编码历史观测，以一维卷积编码未来目标条件，再将两个条件向量拼接后送入MLP，一次性输出完整动作块。BC不包含显式循环状态，因此它主要检验目标中心表示和动作块接口本身是否足以完成动态抓取。

BC-RNN在BC的基础上加入循环时间结构。本文使用GRU而非LSTM作为循环单元\citep{cho2014learning,chung2014empirical}：GRU编码器读取经线性投影的历史观测序列，一维卷积编码未来条件，二者共同初始化GRU解码器。解码器以固定动作查询展开32步动作块，而非训练时读取专家上一动作、推理时读取自身上一预测的自回归结构。该设计避免了训练与推理之间动作输入分布的不一致，同时保持固定长度动作块接口；每16步重新规划时，新的闭环观测反馈自然进入策略。

\subsection{条件变分动作块策略：ACT}

ACT基于动作分块变换器的动作查询解码结构，以条件变分自编码器（CVAE）的隐变量建模动作块条件分布\citep{zhao2023learning,kingma2014autoencoding,sohn2015learning}。模型将历史观测和未来条件投影为带类型嵌入的序列标记，经变换器编码器形成上下文记忆；32个学习动作查询经变换器解码器一次性生成动作块。

训练时，ACT的后验编码器同时读取上下文记忆和专家动作块，输出隐变量分布$q_\phi(z|A^\ast,O,C)$的均值和方差，并经重参数化采样得到$z$。解码器接受该隐变量作为条件，生成动作块并同时预测闭合事件和阶段标签。推理时专家动作块不可用，模型取先验均值$z=0$，因此闭环执行仍是一次确定性前向传播。该结构的作用不是改变输入输出接口，而是在同一动作块接口下检验ACT的变换器条件建模和隐变量动作块分布是否能改善短时轨迹组织、阶段语义和闭合时机。

\subsection{生成式动作块策略：DP}

DP将动作块生成建模为条件去噪问题\citep{ho2020denoising,song2021denoising,chi2023diffusion}。给定专家动作块$A^\ast$，训练时先采样扩散步$s$和高斯噪声$\epsilon$，构造带噪动作块$A_s$；条件1D U-Net\citep{ronneberger2015unet}接收$A_s$、扩散时间嵌入以及从历史观测和未来条件导出的条件向量，预测噪声$\epsilon$。本文DP模型采用FiLM式条件残差块和U-Net下采样/上采样结构，并在去噪主干之外保留闭合和阶段辅助预测头。

推理时，DP从高斯噪声动作块开始，以10步DDIM确定性采样逐步生成32步动作块。相较于BC、BC-RNN和ACT的一次前向传播，DP需要多次去噪网络调用；但由于系统每16步才重新规划一次，其计算开销可被动作块执行窗口分摊。除主模型DP 12M外，另训练1.5M、6M和18M三个容量变体（见附录），以检验容量提升是否转化为闭环成功率增益。

\subsection{训练目标与检查点选择}

BC、BC-RNN和ACT的损失函数为动作$L_1$损失、闭合事件二分类交叉熵和阶段交叉熵的加权和：
\begin{equation}
    \mathcal{L}_{det}
    = \mathcal{L}_{action}
    + \lambda_c \mathcal{L}_{close}
    + \lambda_\phi \mathcal{L}_{phase}.
    \label{eq:deterministic-policy-loss}
\end{equation}
其中$\lambda_c=5.0$，$\lambda_\phi=1.0$。ACT额外加入KL正则项，KL权重为10.0，并设置10000步预热。DP以噪声预测损失为主损失，同时保留相同权重的闭合和阶段辅助监督：
\begin{equation}
    \mathcal{L}_{dp}
    = \mathbb{E}_{s,\epsilon}\|\epsilon_\theta(A_s,s,O,C)-\epsilon\|_2^2
    + \lambda_c \mathcal{L}_{close}
    + \lambda_\phi \mathcal{L}_{phase}.
    \label{eq:diffusion-policy-loss}
\end{equation}

所有主策略共享统一训练配置：使用AdamW优化器\citep{loshchilov2019decoupled}，学习率$2\times10^{-4}$，权重衰减$1\times10^{-6}$，共训练1500个轮次，bf16自动混合精度，启用TF32，EMA衰减为0.999。BC、BC-RNN和ACT按验证动作损失选择检查点；DP按验证噪声预测损失选择检查点。主实验中每个策略--场景--表示组合均采用固定随机种子；因此第四章报告的Wilson置信区间仅反映评估回合的抽样不确定性，不包含训练随机性的种子间方差。验证损失仅用于训练选择与优化诊断，第四章所有主要结论均以闭环成功率为准。

\section{闭环执行与跟踪执行器设置}

闭环执行时，学习策略每16个仿真步重新预测一个32步任务空间动作块。状态机单调推进，不允许阶段回退；闭合预测头可以触发从接近或对齐转入闭合，阶段预测头每步最多推进一个阶段。闭合阈值为0.5，滞回步数为5；转入保持或搬运后，夹爪保持闭合。若系统转入搬运或在保持中稳定10步，则控制切换为搬运控制。

主实验采用MPC跟踪执行器。学习策略输出的任务空间参考在抓取阶段由MPC跟踪器执行，MPC预测时域为10，位置、姿态和平滑权重分别为$w_{pos}=300$、$w_{orient}=50$、$w_{smooth}=5$。若MPC求解失败或非搬运阶段跟踪误差过大，系统以IK风格目标作为回退方案。该设置保持了专家和学习策略之间的执行结构一致性，使策略比较主要反映动作块生成质量。


\section{训练与评估标准}

本文考察两类训练设置：无噪声训练和在线观测噪声增强训练。无噪声训练策略训练时不添加观测噪声；在线观测噪声增强训练策略在训练时注入与噪声评估匹配的观测扰动，包括位置噪声0.03、姿态噪声$10\,\mathrm{deg}$、速度噪声0.02、角速度噪声$5\,\mathrm{deg}$、随机丢弃率0.05、漂移标准差0.01、离群点概率0.02、离群点幅度0.15和2帧时延。噪声评估环境叠加漂移、高斯噪声、离群点、随机丢弃和时延的组合扰动；随机姿态翻滚平移噪声场景通过独立噪声预设构造，不基于无噪观测设置手工修改。图~\ref{fig:ch3-observation-noise-trajectory}以一条真实示教轨迹为例，可视化该噪声栈对立方体位置观测的影响。

\begin{figure}[!htbp]
    \centering
    \includegraphics[width=\textwidth]{fig_ch3_observation_noise_trajectory}
    \caption{无噪观测轨迹与噪声观测轨迹对比。左图在三维空间中叠加同一条立方体位置轨迹的无噪观测和噪声观测，黑色圆点和星号分别标记无噪轨迹的起点和终点，箭头表示运动方向，灰色连线表示部分时刻的观测偏移；线框立方体的一个角点位于坐标原点，其边长为5 cm，用于对比观测噪声与物体尺度。右图给出噪声观测相对无噪观测的位置误差随时间的变化，并标出离群点抽样和随机丢帧导致的陈旧观测。该图仅展示位置分量；训练和评估中的姿态、速度与角速度同样按本文噪声设置受到扰动。}
    \label{fig:ch3-observation-noise-trajectory}
\end{figure}

主评估使用750个闭环评估回合，并报告闭环成功率。记总回合数为$n$，成功回合数为$n_s$，则成功率点估计为$\hat{p}=n_s/n$。本文以95\% Wilson置信区间量化成功率的不确定性\citep{wilson1927probable}：
\begin{equation}
    \mathrm{CI}_\mathrm{Wilson} =
    \frac{\hat{p} + \frac{z^2}{2n} \pm z\sqrt{\frac{\hat{p}(1-\hat{p})}{n} + \frac{z^2}{4n^2}}}{1 + \frac{z^2}{n}},
    \quad z_{0.975}=1.96.
    \label{eq:wilson-ci}
\end{equation}
与Wald区间相比，Wilson区间在$\hat{p}$接近0或1时不会产生零宽度区间，因此即使策略成功率接近0\%或100\%，仍能给出合理的置信范围。750回合评估中，50\%成功率处的区间半宽约为$\pm3.6$个百分点，足以在多数实验条件间区分有意义的性能差异。300回合探索性评估的区间半宽约为$\pm5.7$个百分点，解读相应结果时需考虑这一较大的不确定性。

第四章的噪声鲁棒性分析涉及两类卡尔曼滤波实验：专家FSM-MPC控制器的强噪声评估，以及学习策略在滤波观测下的探索性评估。两类实验共用同一卡尔曼滤波器配置。滤波器作为感知前处理模块插入观测流：观测受到噪声损坏后，先经滤波器完成状态估计，再送入控制器或策略。

滤波器对立方体位置和线速度施加常速度过程模型的卡尔曼滤波\citep{kalman1960new}，过程噪声协方差参数为0.01。位置和速度的测量噪声参数分别设为0.03和0.02，与高斯抖动噪声的标准差一致。对于立方体姿态，滤波器采用指数平滑而非完整的卡尔曼滤波，平滑系数$\alpha=0.3$。该设计是工程上的折中：位置和速度的状态维度较低且动力学为线性，适合卡尔曼滤波；姿态的四元数流形上标准卡尔曼滤波需要更复杂的处理，指数平滑是轻量但有效的替代方案。

滤波后的观测替换原始噪声观测，但不改变下游策略或控制器的接口——滤波器和噪声注入模块都是感知层的组件，学习策略和专家控制器看到的始终是相同格式的观测向量。该滤波器配置同样应用于专家FSM-MPC的强噪声评估（第四章），以保证两类实验之间的可比性。

除主实验外，本文还通过受控消融研究各设计因素的影响。表示消融比较完整目标中心表示、紧凑目标中心表示、完整世界坐标表示和紧凑世界坐标表示；历史长度消融比较紧凑目标中心表示的不同$H_o$；数据量消融比较10\%、25\%、50\%和100\%数据；容量实验比较BC-RNN和DP扩容结果。所有验证损失仅用于训练选择与优化诊断，不能替代闭环性能。
